\chapter{Voorwoord}

Dit boek heeft een eenvoudige ambitie: één samenhangend boek over aard- en
ruimtewetenschappen dat alles bestrijkt van het eerste schooljaar tot het
einde van de bachelor, geschreven in het Nederlands, voor lezers waar ook
ter wereld.

De stijl is beknopt en streng: lesteksten opgebouwd uit definities,
voorbeelden, proposities, methoden en schema's, met zorgvuldige
redeneringen zodra die op het betreffende niveau toegankelijk zijn,
gevolgd door zorgvuldig gekozen oefeningen. Resultaten die zonder
volledige afleiding worden gegeven, zijn uitdrukkelijk als aangenomen
gemarkeerd. Van alle oefeningen staan achter in het boek volledige
uitwerkingen. Elk getal dat over de echte Aarde, haar oceanen, haar
atmosfeer of de hemel wordt genoemd, komt uit een naslagbron, en elke
kromme die op echte metingen berust, is berekend uit de gepubliceerde
gegevens.

De cursus bestaat uit vijf boeken. Het eerste bestrijkt de jaren van de
basisschool en de onderbouw, het tweede de drie jaren van de bovenbouw, en
de laatste drie de drie universiteitsjaren. De twee schoolboeken leggen
nooit twee keer hetzelfde begrip uit: elk idee wordt één keer behandeld,
in het jaar waarin het voor het eerst past, en een later hoofdstuk dat het
nodig heeft, begint met een korte herinnering, \emph{Wat je al weet},
voordat het verder gaat. Waar een begrip bij de natuurkunde, de biologie of
de scheikunde hoort, leggen de zusterboeken uit dezelfde reeks het uit, en
dit boek brengt het in herinnering voordat het de blik van de
aardwetenschapper aanneemt. In de schoolboeken is elk jaar een deel,
verdeeld in hoofdstukken; elk universiteitsboek is één jaar, verdeeld in
hoofdstukken. Hoe vroeger het jaar, hoe jonger de lezer voor wie het
geschreven is: meer illustraties, meer uitgewerkte tussenstappen en
mildere oefeningen.

\section*{Hoe je dit boek leest}

De uitspraken worden per hoofdstuk genummerd en zijn met kleur gecodeerd:
definities in blauw, stellingen in rood, proposities en verwanten in
oranje, en methoden --- korte recepten voor standaardtaken --- in groen.
De oefeningen lopen van ($\star$) voor rechtstreekse toepassingen tot
($\star\star\star$) voor lastigere opgaven. Omkaderde blokken bevatten
wat niet bij de redenering zelf hoort: een herinnering aan wat een eerder
hoofdstuk heeft uitgelegd, hoe geologen en andere aardwetenschappers in het
veld of in het laboratorium werken, een bladzijde geschiedenis, cijfers die
van jaar tot jaar veranderen (met het jaar waarvoor ze gelden), en
veiligheidsadvies.

\section*{Veiligheid}

De Aarde leer je kennen door ernaar te kijken: naar de hemel, het weer, een
landschap, de stenen van een gebouw, de bodem van een tuin. Meer vraagt dit
boek nooit van zijn lezer. Kijk nooit naar de Zon, met of zonder
instrument: je wordt er blind van. Groeven, kliffen, grotten, rivieren die
buiten hun oevers treden en vulkanen zijn plaatsen voor geoefende
professionals met de juiste uitrusting, en dit boek beschrijft hun werk
zonder iemand erheen te sturen.

\section*{Meewerken}

De bron van dit boek staat in een openbare git-repository:
\begin{center}
\url{https://github.com/bvirrion/one-earth-and-space-sciences-book}
\end{center}
Bijdragen --- nieuwe hoofdstukken, correcties, betere uitleg, extra
oefeningen --- zijn welkom; zie \texttt{CONTRIBUTING.md} in de hoofdmap van
de repository.
